\section{Method and reproducible commands}

This edition used source reads, Git comparisons, publication checks and XeLaTeX. It ran no
Cargo or Clippy command, VM, real-GPU gate, provider operation or remote-host experiment.
Every hardware result is carried from a dated source. Publication checks here do not turn
those records into measurements of the cover revision.

\subsection{Sources and code distance}

\code{sources_v3.json} records repository revisions, paths and Git blobs, plus the three
external primary sources checked on the publication date. Numeric citations link directly
to those sources. The full mapping, historical evidence names, PR metadata and gap topics are
in the \href{\repobase/blob/codex/whitepaper-v3-20261027/docs/whitepaper/kayfabe_v3_index.md}
{companion index, \code{kayfabe_v3_index.md}}. That branch link moves with publication;
the source revisions recorded inside it do not. External pages can change after checking.

The following reconstructs repository-source identity and the result-distance calculation:
\begin{verbatim}
python3 - <<'PY'
import json, subprocess
from pathlib import Path
p = Path('docs/whitepaper/sources_v3.json')
for row in json.loads(p.read_text())['sources']:
    ref = row['revision'] + ':' + row['path']
    blob = subprocess.check_output(['git', 'rev-parse', ref], text=True)
    assert blob.strip() == row['blob'], ref
PY
# PIN is the cover revision; RESULT is a runtime from the results table.
git diff --name-only RESULT PIN -- crates qemu cuda
git rev-list --left-right --count PIN...RESULT
\end{verbatim}

Changed paths include additions, removals and modifications. Commit counts cover the whole
Git graph, not only the three product directories. Zero result-only commits means ancestor;
it is not a passing-test status. A divergent result has commits absent from the pin.
\code{result_distances_v3.json} preserves full runtime hashes and rerun status. No result in
that file was rerun on hardware for this edition.

\subsection{PR snapshot and historical gap index}

\begin{verbatim}
gh pr list --state open --limit 200 \
  --json number,title,headRefName,headRefOid,baseRefName,isDraft,url
gh pr list --state all --limit 250 \
  --json number,headRefName,baseRefName,state,url
\end{verbatim}

The snapshot was refreshed after the revision work. It freezes each open PR's number, base,
head, draft state and known parent PR: a parent is identified by an exact base-branch to
PR-head-branch match, including merged parents. A branch base without such a match remains
unknown. A recorded parent is metadata, not proof that the full dependency was reviewed.
Live heads, draft states and bases will change after capture.

The companion also preserves 101 historical gap topics. They include closed and superseded
rows; the count is not a current defect count. Application gaps, capability gaps and postponed
features are identified separately there. No template inventory occupies the main argument.

\subsection{Build and verify the publication}

\begin{verbatim}
python3 -m unittest discover -s docs/whitepaper -p test_build_v3.py
python3 scripts/ci/debt.py claims
git diff --check
# Commit TeX inputs before producing a publishable PDF.
bash docs/whitepaper/build.sh v3
python3 docs/whitepaper/build_v3.py --check
pdfinfo docs/whitepaper/kayfabe_v3.pdf
pdftotext -layout docs/whitepaper/kayfabe_v3.pdf -
\end{verbatim}

\readcode The v3 target runs XeLaTeX twice with halt-on-error and an input recorder. It checks
for TeX errors, missing glyphs and undefined references. Recorded repository inputs must match
committed blobs. A JSON receipt binds input hashes, source commit, compiler and output digest.
A later PDF-only commit leaves the input blobs unchanged. This is a local publication check,
not a signed attestation of a trusted compiler or host.

\subsection{Reproducing hardware evidence is separate work}

The following are repository entry points, not commands run for this paper. They require an
authorised, exclusively leased test machine and configuration capture.\src{rules}
\begin{verbatim}
bash scripts/bench/v3_gates.sh
KF_DEVICE=kf3 bash scripts/fastguest/fast_suite.sh <tag> 180
bash scripts/bench/cuda_ladder.sh host <tag>
bash scripts/bench/cuda_ladder.sh guest <tag>
bash scripts/bench/perf_compare.sh <rev> \
  --baseline <rev-or-result-dir> --lanes bare,fast --reps 3
\end{verbatim}

\doc The performance method requires source revision, actual die, drivers, nested status,
placement, settings, flags, sample unit, failed runs and the control's spread. Changing a paper's
source pin does not update a hardware receipt.\src{linuxperf}

For the stall case, preregister completion, recovery and latency outcomes separately. Keep
default-watchdog and longer-observation arms. Missing scheduler preemption timestamps,
packet association, native GPU clocks or guest events remain missing observations. A capped
trace ring cannot silently shrink a failure denominator.\src{latency}\src{msrecovery}

\section{Glossary and historical names}

\begin{description}[style=nextline,itemsep=0.25em]
\item[ABI] Application binary interface: layouts and calling conventions.
\item[BAR] PCI address window. Only BAR0 writes trap in this design.
\item[CE / GR] Copy engine / graphics-compute engine in the source records.
\item[GPFIFO / GP / pushbuffer] Ring of command references / one reference / command bytes.
\item[GET / PUT] Fetch cursor / publication cursor; fetching does not establish execution.
\item[IOVA / GPA] Device-visible I/O virtual address / guest physical address.
\item[IRQ / MSI / DPC] Interrupt request / message-signalled interrupt / deferred Windows work.
\item[MEM\_OP / split] Memory-operation command / ordering boundary for a required map refresh.
\item[PTE / PDE / PDB] Page-table entry / page-directory entry / page-directory base.
\item[RM / GSP] Resource manager / GPU System Processor.
\item[TSG / USERD] Timeslice group / channel producer-consumer state memory.
\item[USER] GPU channel privilege; distinct from host request permission and Unix credentials.
\item[UVM] Unified virtual memory, including fault delivery and replay.
\item[WDDM / TDR] Windows Display Driver Model / timeout detection and recovery.
\item[ogkm] NVIDIA's open GPU kernel modules, used as pinned reference sources.
\end{description}

\subsection{Decoding old file names}

These are lookup aids for historical sources, not extra concepts required by the argument.
\begin{description}[style=nextline,itemsep=0.25em]
\item[v3 / pre-v3] Current \code{kf-*} product / frozen \code{kayfabe-*} predecessor.
\item[M / B-prime / baton] Alternative relay-handoff designs.\src{alternativerelay}
\item[ENS5 / ENS6 / ENS7] Sparse-admission undo / host-call churn and slices /
framebuffer-scaled sparse budgets.\src{admission}\src{rmcost}\src{budget}
\item[ENS8 / ENS10 / ENS11] Consumed-ring wakes / mapping scope / split-latency continuation.
\src{owed}\src{walknarrow}\src{latency}
\item[T-space] Address space reserved for privileged rewritten work.\src{rulings}
\item[Boot labels and candidate numbers] Harness identifiers, not release versions or grades.
The companion records the few identifiers needed to find the unverified analysis.
\end{description}

\section{Change record}

The predecessor describes the earlier multi-process architecture and remains superseded.
The first v3 edition replaced that model, pinned the source, introduced explicit limitations
and added the committed-input publication gate.\src{old}\src{architecture}

This cold-reader revision states the stall hypothesis before the internals, distinguishes
CODE from DOC and REPORTED from UNVERIFIED, restores positive Windows results and native
performance denominators, and makes host-plus-guest performance an explicit missing result.
It replaces generic concept boxes with four mechanism diagrams plus the lifecycle diagram.
The full inventories move to the companion; 36 questions now state their own experimental
setup and decision. Neither this revision nor its Git-distance table qualifies new hardware.
